% Generated draft: requires mathematical and source review.
\documentclass[11pt]{article}
\usepackage[margin=25mm]{geometry}
\usepackage{fontspec}
\setmainfont{Latin Modern Roman}
\usepackage{amsmath,amssymb,mathtools,unicode-math}
\setmathfont{Latin Modern Math}
\usepackage{enumitem,needspace,xurl,hyperref}
\setlength{\parindent}{0pt}
\setlength{\parskip}{5pt}
\setlength{\emergencystretch}{3em}
\newcommand{\N}{\mathbb{N}}
\newcommand{\Z}{\mathbb{Z}}
\newcommand{\Q}{\mathbb{Q}}
\newcommand{\R}{\mathbb{R}}
\newcommand{\C}{\mathbb{C}}
\newcommand{\F}{\mathbb{F}}
\newcommand{\A}{\mathbb{A}}
\newcommand{\PP}{\mathbb{P}}
\newcommand{\E}{\mathbb{E}}
\newcommand{\eps}{\varepsilon}
\newcommand{\dd}{\,\mathrm{d}}
\begin{document}
\section*{UnsolvedMath ID: 701457}
\textbf{Source problem code:} catalogue entry
\subsection{Definitions and mathematical statement}
The source does not specify the ambient class of groups or define a formation. The following makes its question precise relative to an ambient class \(\mathcal U\), using closure under finite subdirect products. Whether this is the intended convention requires review; in particular, the source does not explicitly restrict to finite groups.

Fix a nonempty class \(\mathcal U\) of groups closed under isomorphisms, homomorphic images, and finite subdirect products. A \emph{formation in \(\mathcal U\)} is a nonempty subclass \(\mathcal F\subseteq\mathcal U\) closed under these operations. Here a subdirect product of \(G_1,\ldots,G_n\), with \(n\geq 1\), is a subgroup \(S\leq G_1\times\cdots\times G_n\) whose projection onto each factor is surjective. For \(G\in\mathcal U\), define
\[
\operatorname{form}_{\mathcal U}(G)
=\bigcap\{\mathcal F:\mathcal F\text{ is a formation in }\mathcal U,\ G\in\mathcal F\}.
\]
A formation is \emph{one-generator} if it equals \(\operatorname{form}_{\mathcal U}(G)\) for some group \(G\in\mathcal U\). This refers to generation of a formation by one group and does not require that group to be cyclic. A \emph{subformation} of \(\mathcal F\) is a formation \(\mathcal H\) with \(\mathcal H\subseteq\mathcal F\).

\textbf{Question (V. P. Burichenko, Kourovka Notebook 18.28).} For the intended ambient class and formation convention, is it true that
\[
\forall G\in\mathcal U\quad
\forall\mathcal H\subseteq\operatorname{form}_{\mathcal U}(G),
\qquad
\mathcal H\text{ is a formation}
\ \Longrightarrow\ 
\exists H\in\mathcal U\quad
\mathcal H=\operatorname{form}_{\mathcal U}(H)?
\]
The choice of \(H\) may depend on both \(G\) and \(\mathcal H\). This is a formulation schema pending confirmation of the source's ambient class and closure convention, not an assertion that the question ranges over every possible \(\mathcal U\).
\subsection{Short English statement}
Must every subformation of a formation generated by a single group itself be generated by a single group? The source leaves the ambient class of groups and the formation convention unstated.
\subsection{Sources}
\begin{itemize}
\item Supplied local UnsolvedMath registry record: permanent ID \(701457\), exact source problem code \texttt{catalogue entry}, title \emph{Kourovka Notebook Problem 18.28}; upstream revision \texttt{d371ee917ac9647fbb06cd6fc529deec6f3dda82}.
\item The supplied record attributes the problem to V. P. Burichenko, proposed in 2014, and cites \emph{The Kourovka Notebook (Unsolved Problems in Group Theory)}, 18th issue (2014), problem 18.28, with the 21st issue edition (2026), \texttt{arXiv:1401.0300v48}. Its collection description names E. I. Khukhro and V. D. Mazurov as editors.
\item The registry labels the problem \emph{open}. This is supplied source metadata, not an independently verified assessment. The cited edition and its surrounding definitions have not been independently consulted for this draft.
\end{itemize}
\end{document}
